\documentclass[10pt, aspectratio=169]{beamer}
\input{../../_en_preambulo_beamer}
% Comandos y macros estadísticas compartidas para las presentaciones Beamer en inglés (Metropolis)

% Conjuntos numéricos básicos
\newcommand{\R}{\mathbb{R}}
\newcommand{\N}{\mathbb{N}}
\newcommand{\Q}{\mathbb{Q}}
\newcommand{\Z}{\mathbb{Z}}
\newcommand{\set}[1]{\left\{ #1 \right\}}
\newcommand{\abs}[1]{\left|#1\right|}
\newcommand{\norm}[1]{\left\| #1 \right\|}

% Comandos estadísticos y probabilísticos del curso
\newcommand{\Var}{\operatorname{Var}}
\newcommand{\cov}{\operatorname{Cov}}
\newcommand{\corr}{\rho}
\newcommand{\comb}[2]{\begin{pmatrix} #1 \\ #2 \end{pmatrix}}
\newcommand{\s}{\sigma}
\newcommand{\particion}{\mathfrak{P}}
\newcommand{\card}[1]{\#\left( #1 \right)}
\newcommand{\seq}[3]{\set{#1}_{#2}^{#3}}
\newcommand{\E}{\mathbb{E}}
\newcommand{\Prob}{\mathbb{P}}

% Entornos pedagógicos y cajas en inglés académico natural (soportando tanto nombres en inglés como en español)
\newenvironment{discussion}{%
  \begin{alertblock}{Discussion Question}%
}{%
  \end{alertblock}%
}
\newenvironment{discusion}{%
  \begin{alertblock}{Discussion Question}%
}{%
  \end{alertblock}%
}

\newenvironment{reflection}{%
  \begin{alertblock}{Classroom Reflection}%
}{%
  \end{alertblock}%
}
\newenvironment{reflexion}{%
  \begin{alertblock}{Classroom Reflection}%
}{%
  \end{alertblock}%
}

\newenvironment{paradox}{%
  \begin{alertblock}{Exploratory Paradox}%
}{%
  \end{alertblock}%
}
\newenvironment{paradoja}{%
  \begin{alertblock}{Exploratory Paradox}%
}{%
  \end{alertblock}%
}

\newenvironment{motivation}{%
  \begin{exampleblock}{Motivation: Data Science in Practice}%
}{%
  \end{exampleblock}%
}
\newenvironment{motivacion}{%
  \begin{exampleblock}{Motivation: Data Science in Practice}%
}{%
  \end{exampleblock}%
}

\newenvironment{quickchallenge}{%
  \begin{block}{Quick Team Challenge}%
}{%
  \end{block}%
}
\newenvironment{retoclase}{%
  \begin{block}{Quick Team Challenge}%
}{%
  \end{block}%
}


\title[One-Way ANOVA]{Section 08.01: One-Way Analysis of Variance}
\subtitle{Experimental Design, Variance Decomposition, and the $F$-Test}
\author[J. Castillo Colmenares]{Juliho Castillo Colmenares}
\institute{Tecnológico de Monterrey}
\date{\vspace{-1.5cm}}

\begin{document}

% Slide 1: Title
\begin{frame}[plain]
  \titlepage
\end{frame}

% Slide 2: Roadmap
\begin{frame}{Roadmap --- Chapter 08: Elements of Experimental Design (ANOVA)}
  \footnotesize
  Where are we in the modular development of this chapter? \pause
  \begin{itemize}\itemsep=0.08em
    \item \textbf{08.01} One-Way Analysis of Variance (ANOVA) and the $F$-Test \emph{(Today)} \pause
    \item \textbf{08.02} ANOVA Assumptions, the Levene/Bartlett Tests, and Diagnostics
  \end{itemize}
  \vspace{-0.05cm}
  \begin{block}{Session Objective}
    Generalize the comparison of two means (Chapter 07) to the simultaneous comparison of $k \geq 3$ population means without inflating the Type I error rate: build the one-way ANOVA linear model, decompose total variance, derive the $F$ statistic, and apply post-hoc procedures (LSD, Tukey, Bonferroni) to identify exactly which treatment pairs differ.
  \end{block}
\end{frame}

% Slide 3: Motivation
\begin{frame}{The Problem of Comparing $k$ Groups Simultaneously}
  \footnotesize
  \begin{motivacion}
    Comparing $k=5$ machine learning algorithms by running $\binom{5}{2}=10$ pairwise $t$-tests, each at $\alpha=0.05$, produces a family-wise Type I error rate of $1-(0.95)^{10}\approx 40.13\%$ ---even if all five algorithms are in truth identical.
  \end{motivacion}

  \vspace{0.15cm}
  \pause
  \begin{itemize}\itemsep=0.1cm
    \item \textbf{Fisher's solution:} a single global test ($F$) that contrasts all $k$ means simultaneously. \pause
    \item \textbf{Key idea:} partition the total observed variance into a component \emph{between} treatments and a component \emph{within} each treatment (error).
  \end{itemize}
\end{frame}

% Slide 4: Model and Decomposition
\begin{frame}{Linear Model and Sum-of-Squares Decomposition}
  \footnotesize
  \begin{block}{Additive Parametric Model}
    $$Y_{ij} = \mu + \tau_i + \epsilon_{ij}, \qquad \epsilon_{ij}\sim N(0,\sigma^2) \text{ i.i.d.}, \qquad \sum_{i=1}^k n_i\tau_i = 0$$
  \end{block}
  \pause

  \vspace{0.1cm}
  \begin{alertblock}{Partition Theorem}
    $$\underbrace{\sum_i\sum_j (Y_{ij}-\bar Y_{..})^2}_{\text{SST},\ \nu_T=N-1} = \underbrace{\sum_i n_i(\bar Y_{i.}-\bar Y_{..})^2}_{\text{SSTR},\ \nu_{TR}=k-1} + \underbrace{\sum_i\sum_j (Y_{ij}-\bar Y_{i.})^2}_{\text{SSE},\ \nu_E=N-k}$$
  \end{alertblock}
\end{frame}

% Slide 5: ANOVA Table and F Statistic
\begin{frame}{Mean Squares and the $F$ Test Statistic}
  \footnotesize
  \begin{block}{Expected Values}
    $$E(\text{MSE})=\sigma^2, \qquad E(\text{MSTR}) = \sigma^2 + \frac{\sum_i n_i\tau_i^2}{k-1}$$
  \end{block}
  \pause

  \vspace{0.1cm}
  \begin{alertblock}{Decision Rule}
    $$F = \frac{\text{MSTR}}{\text{MSE}} = \frac{\text{SSTR}/(k-1)}{\text{SSE}/(N-k)} \sim F_{k-1,\,N-k} \quad\text{under } H_0; \qquad \text{Reject } H_0 \text{ if } F>F_{\alpha,k-1,N-k}$$
  \end{alertblock}
\end{frame}

% Slide 6: Post-hoc
\begin{frame}{Post-Hoc Procedures: Which Pairs Differ?}
  \footnotesize
  If $F$ is significant, we know at least one pair differs, but not which. \pause
  \begin{itemize}\itemsep=0.1cm
    \item \textbf{Fisher's LSD:} $|\bar Y_{i.}-\bar Y_{j.}| > t_{\alpha/2,N-k}\sqrt{\text{MSE}(1/n_i+1/n_j)}$. Powerful but does not control the FWER. \pause
    \item \textbf{Tukey's HSD:} uses the studentized range distribution $q$; controls the FWER exactly for all pairwise comparisons. \pause
    \item \textbf{Bonferroni Correction:} adjusts $\alpha^* = \alpha/\binom{k}{2}$; universal but conservative.
  \end{itemize}
\end{frame}

% Slide 7: Applications
\begin{frame}{Applications in Data Science and Engineering}
  \footnotesize
  \begin{itemize}\itemsep=0.1cm
    \item \textbf{Algorithm benchmarking:} comparing $F_1$-score across multiple classification architectures. \pause
    \item \textbf{A/B/n testing:} evaluating several campaign or interface variants simultaneously. \pause
    \item \textbf{Systems engineering:} comparing latency across server configurations or protocols. \pause
    \item \textbf{Randomized Block Design (RCBD):} when experimental units are not homogeneous, blocking known noise dramatically increases test power.
  \end{itemize}
\end{frame}

% Slide 8: Python Lab (Block 1)
\begin{frame}[fragile]{Python Lab: One-Way ANOVA and the $F$-Test}
  \scriptsize
  SST = SSTR + SSE decomposition for 4 server configurations, verified against \texttt{scipy.stats.f\_oneway} (\texttt{08.01\_one\_way\_anova.py}):
  {\lstset{basicstyle=\fontsize{3.6pt}{4.3pt}\selectfont\ttfamily,aboveskip=0.05em,belowskip=0.05em}
  \lstinputlisting[firstline=17, lastline=39]{../../code/08_diseno_experimentos/08.01_one_way_anova.py}}
\end{frame}

% Slide 9: Python Lab (Block 2)
\begin{frame}[fragile]{Python Lab: Tukey's HSD}
  \scriptsize
  Pairwise multiple comparisons using the studentized range from \texttt{scipy.stats.studentized\_range}:
  {\lstset{basicstyle=\fontsize{3.8pt}{4.5pt}\selectfont\ttfamily,aboveskip=0.05em,belowskip=0.05em}
  \lstinputlisting[firstline=48, lastline=63]{../../code/08_diseno_experimentos/08.01_one_way_anova.py}}
\end{frame}

% Slide 10: Python Lab (Block 3)
\begin{frame}[fragile]{Python Lab: Randomized Complete Block Design}
  \scriptsize
  SST = SSTR + SSB + SSE decomposition and the Relative Efficiency of blocking:
  {\lstset{basicstyle=\fontsize{3.4pt}{4.1pt}\selectfont\ttfamily,aboveskip=0.05em,belowskip=0.05em}
  \lstinputlisting[firstline=66, lastline=92]{../../code/08_diseno_experimentos/08.01_one_way_anova.py}}
\end{frame}

% Slide 11: Terminal Output
\begin{frame}[fragile]{Python Lab: Terminal Output}
  \scriptsize
  Standard output produced when running the lab script:
  \vspace{0.02cm}
  \begin{block}{Standard Console Output}
    \fontsize{3.8pt}{4.6pt}\selectfont
    \begin{verbatim}
=== Block 1: One-Way ANOVA ===
SSTR=3754.20 (df=3), SSE=979.26 (df=20), F=25.5580, p=0.000000
scipy.stats.f_oneway check: F=25.5580, p=0.000000

=== Block 2: Tukey HSD ===
q(0.05,4,20)=3.9583, HSD=11.3075
Config B vs Config C: 7.52 -> not significant (all other pairs SIGNIFICANT)

=== Block 3: RCBD ===
SST=437.67 = SSTR(52.67) + SSB(382.00) + SSE(3.00)
F_TR=52.67 (p=0.000157), F_B=254.67 (p=0.000001), Relative Efficiency=70.18
    \end{verbatim}
  \end{block}
\end{frame}


% Slide 16: Closing
\begin{frame}{Section Closing: One-Way ANOVA}
  \begin{reflexion}
    We have generalized the comparison of two means to $k\geq3$ groups without inflating the Type I error rate: ANOVA partitions total variance into a between-treatments component and an error component, and the $F$ statistic formalizes the decision. Post-hoc procedures (LSD, Tukey, Bonferroni) complete the analysis by identifying exactly which pairs differ.
  \end{reflexion}

  \vspace{0.2cm}
  \pause
  \begin{block}{Key Takeaways}
    \begin{itemize}\itemsep=0.08cm
      \item \textbf{Decomposition:} $\text{SST}=\text{SSTR}+\text{SSE}$, the foundation of all experimental design.
      \item \textbf{$F$ statistic:} the ratio of two independent estimators of $\sigma^2$ under $H_0$.
      \item \textbf{Blocking (RCBD):} incorporating known noise as a separate factor dramatically increases power.
    \end{itemize}
  \end{block}
\end{frame}

% Slide 17: Modular Perspective
\begin{frame}{Modular Perspective --- The Next Challenge}
  \scriptsize
  \begin{retoclase}
    Everything demonstrated today --- the $F$ distribution of the test statistic, the validity of the post-hoc procedures --- rests on three assumptions of the parametric model: normality, homoscedasticity, and independence of the errors. Section 08.02 (chapter closing) rigorously audits these assumptions with the Levene, Bartlett, and Shapiro-Wilk tests, and presents the nonparametric Kruskal-Wallis alternative when they fail.
  \end{retoclase}

  \vspace{0.08cm}
  \pause
  \begin{alertblock}{Recommended Preparation}
    \begin{itemize}\itemsep=0.06cm
      \item Review why ANOVA is robust to homoscedasticity violations only when sample sizes are equal.
      \item Reflect on what happens to the validity of the $F$-test if the errors are not independent.
    \end{itemize}
  \end{alertblock}
\end{frame}

\end{document}
